\section{三阶并列挖角构型的四阶损失}\label{four:section}

定理~\ref{cut:main}中 $k=1$ 与 $k=0$ 具有相同的三阶系数。
本节进一步分辨这两个构型。四阶传递只需几乎处处的径向二阶展开，
而一致支配由上一节已经证明的余项提供。

\begin{theorem}[单位挖角比例的四阶项]\label{four:main}
对 $E_{1,t}$ 有
\begin{equation}\label{four:expansion}
 \frac{|H_{1,\mathcal Q_c}(1-t,E_{1,t})|}{|E_{1,t}|}
 =2t^2+2t^3-\frac23t^4+o(t^4).
\end{equation}
因此正方形输入的面积比与此面积比之差为
$\frac83t^4+o(t^4)$；矩形测试上确界 $R(1-t)$ 与此面积比之差为
$\frac{19}{6}t^4+o(t^4)$。两个比较中的四阶损失都不可提高为 $o(t^4)$。
\end{theorem}
\begin{proof}
由对称性，令射线方向为 $(1-s,s)$，固定 $0<s<1/2$。
上一节给出径向边界
\[
 r_t(s)=1+2st+O(t^2),
\]
其中余项对整个 $[0,1/2]$ 一致。特别地，对固定正 $s$，小 $t$ 时 $r_t(s)>1$。
记 $r=1+2st+qt^2$，$q$ 在有界区间内变化，并置
$v_1=r(1-s)$、$v_2=rs$。

由命题~\ref{rect:union-scales}，所有尺度候选可写为 $\ell=1+ta$，其中
\[
 a\in\{\pm2v_1,\ \pm2v_2,\ -2\pm2v_1,\ -2\pm2v_2\}.
\]
当 $r=1$ 时，平坦尺度带为 $[-2s,2s]$，且二阶目标在该带上恰等于 $-a$。
故该带上的唯一最佳尺度为 $a=-2s$。
对固定 $0<s<1/2$，带外的候选要么有严格的一阶损失，
要么在带内有严格的二阶损失。唯一在 $r\to1$ 时与 $-2s$ 重合的两个候选是
\[
 a_e=-2v_2,\qquad a_c=-2+2v_1.
\]
其余六个候选的间隙可逐一列出（以下取 $r=1$）：
\[
\begin{array}{c|c|c}
 a\text{ 的候选}&D-1&2s-F\text{（仅在 }D=1\text{ 时）}\\\hline
 \pm2v_1&1-2s&--\\
 2v_2&0&4s\\
 -2-2v_1&3-2s&--\\
 -2-2v_2&1+2s&--\\
 -2+2v_2&1-2s&--
\end{array}
\]
所以这些间隙对每个固定内角均严格为正，在充分小 $t$ 时不影响边界附近的最大值。
这里的起效阈值可以依赖 $s$；最后的面积传递使用一致支配，而不要求它对 $s$ 一致。

对上述固定方向和充分小 $t$，$a_e$ 的精确归一化平均为
\[
 V_e=\frac{(1-tr)(1-2trs)-t^2(2-2rs)}{(1-2trs)\sqrt{1-4t^2}}.
\]
在 $a_c$ 处，令 $b_1=v_1-a_c/2$、$b_2=v_2-a_c/2$，则平均为
\[
 V_c=\frac{(1-tb_1)(1-tb_2)-t^2(2-v_2+a_c/2)}
 {(1+ta_c)\sqrt{1-4t^2}}.
\]
这些分支成立的原因是 $r>1$ 且 $0<s<1/2$ 固定；角块的长坐标负侧捕获为零，
短坐标负侧捕获保持为正。直接展开，关于有界 $q$ 一致地有
\begin{align*}
 V_e-(1-t)&=t^3[-2-4s+8s^2-q]+O_s(t^4),\\
 V_c-(1-t)&=t^3[-2-4s+4s^2-q]+O_s(t^4).
\end{align*}
两者之差的首项为 $4s^2t^3>0$，故边界附近真正的极大值由 $a_e$ 提供。
分别取 $q=Q(s)\pm\varepsilon$ 作严格高低夹逼，得到
\begin{equation}\label{four:radial-limit}
 \frac{r_t(s)-1-2st}{t^2}\longrightarrow
 Q(s):=-2-4s+8s^2\qquad(0<s<1/2).
\end{equation}

定理~\ref{cut:main}的统一径向余项保证式~\eqref{four:radial-limit}左侧一致有界。
因而可对 $s$ 使用支配收敛；两个端点无需另行计算，其测度为零。
面积恒等式及质量归一化给出四阶系数
\[
 8\int_0^{1/2}Q(s)\,ds+
 4\int_0^{1/2}(2s)^2\,ds+8
 =-\frac23.
\]
这里最后的 $8$ 来自除以 $|E_{1,t}|=1-4t^2$。
因此得到式~\eqref{four:expansion}。
正方形面积比为 $2t^2/(1-t)=2t^2+2t^3+2t^4+O(t^5)$，
而 $R(1-t)=2t^2+2t^3+\frac52t^4+O(t^5)$，相减即得两个比较及其sharp阶。
\end{proof}

\begin{remark}\label{four:scope}
上面的 $o(t^4)$ 不自动给出五阶展开；端点方向的移动分支可能影响更高阶项。
四阶损失也只比较这三个指定测试对象，不能据此认定矩形类给出全体最优高阶系数。
\end{remark}
